\chapter{Conclusão}
\label{sec:conclusao}

Este trabalho investigou a reconstrução virtual de temperaturas
a partir de medições de sensores IoT instalados no ambiente 1006
do LASDPC. O objetivo proposto foi atendido por meio da preparação,
sincronização e enriquecimento de dados históricos de nove sensores
físicos, cujas medições foram agregadas em intervalos de cinco
minutos. A avaliação permitiu comparar métodos de reconstrução
e investigar a influência da quantidade e da seleção de sensores
de referência, bem como a contribuição de informações ambientais
complementares. A principal contribuição consiste na avaliação
experimental dessas estratégias em uma infraestrutura real
preexistente, com procedimentos e resultados documentados
para reprodução.

Na comparação principal, os métodos supervisionados apresentaram
erros médios inferiores aos da média aritmética. A regressão
linear obteve MAE médio entre sensores-alvo de 0,257~°C,
enquanto a Random Forest alcançou 0,254~°C, em comparação
a 0,927~°C da referência. A Random Forest apresentou o menor
MAE médio, enquanto a regressão linear obteve a menor média
dos RMSE. Entretanto, as diferenças agregadas entre os modelos
supervisionados foram pequenas e o desempenho variou conforme
o sensor-alvo. A média aritmética, por exemplo, apresentou
menor MAE nos sensores 1 e 86782, evidenciando que nenhuma
das abordagens foi superior em todos os casos.

Diante desses resultados, a regressão linear se apresenta
como uma alternativa inicial competitiva para a reconstrução
de temperaturas no cenário estudado, considerando sua
simplicidade estrutural e seu desempenho próximo ao da
Random Forest. Essa interpretação não estabelece superioridade
estatística nem exclui a utilização de outros métodos,
especialmente em sensores para os quais foram observadas
vantagens individuais.

A investigação da quantidade de entradas mostrou que
utilizar todos os sensores disponíveis não necessariamente
produz estimativas mais precisas. Na regressão linear com
seleção por correlação calculada exclusivamente no treinamento,
quatro sensores de referência apresentaram o menor MAE médio
entre as configurações avaliadas, de 0,238~°C. Esse resultado
sugere que subconjuntos de sensores podem oferecer informações
suficientes para reconstruções competitivas. Entretanto,
o desempenho depende do sensor-alvo e da regra de seleção
empregada, não sendo possível estabelecer uma quantidade
universalmente ótima com base nos experimentos realizados.

Na análise espacial complementar, realizada com seis sensores
de posições conhecidas, a seleção por correlação histórica
apresentou menor MAE médio que a seleção por proximidade
física para uma a quatro entradas. Esse comportamento indica
que a distância euclidiana, isoladamente, não foi o critério
mais favorável nas configurações avaliadas. A representação
espacial das temperaturas médias também permitiu visualizar
diferenças entre os dispositivos, utilizando interpolação
IDW exclusivamente para fins ilustrativos, sem validar
um campo térmico contínuo.

O enriquecimento das entradas com umidade relativa produziu
uma redução pequena no MAE médio da regressão linear, com
melhorias em cinco sensores-alvo e aumento do erro nos
outros quatro. Dessa forma, não foi observada uma vantagem
uniforme da inclusão dessa grandeza nas condições avaliadas.
Os registros de energia, por sua vez, apresentaram baixa
disponibilidade temporal, cobrindo apenas 4,07\% da grade
considerada, e não tiveram sua unidade física confirmada.
Por esse motivo, sua utilização permaneceu restrita a uma
análise exploratória.

As conclusões obtidas estão limitadas a um único ambiente,
um período de observação e uma divisão cronológica de
treinamento e teste, considerando apenas janelas com
temperaturas disponíveis para todos os sensores do conjunto
principal. Não foram avaliados cenários de indisponibilidade
simultânea das entradas ou a estabilidade dos modelos em
outras condições ambientais. Além disso, as medições dos
próprios sensores foram utilizadas como referência,
sem comparação com instrumentos calibrados. Portanto,
os erros observados não demonstram exatidão metrológica
nem validam diretamente a reconstrução de temperaturas
em posições não instrumentadas ou a avaliação de conforto
térmico.

Como continuidade, recomenda-se avaliar os métodos em
outros períodos e ambientes, investigar estratégias de
reconstrução diante de falhas dos sensores de referência
e validar estimativas em posições não instrumentadas.
A integração de novas fontes de dados, após a verificação
de suas convenções temporais e unidades físicas, também
poderá ampliar as possibilidades de análise. Por fim,
a utilização das temperaturas reconstruídas em estudos
de conforto térmico dependerá da definição de um protocolo
específico e da obtenção das demais variáveis ambientais
e pessoais necessárias.

Em síntese, os resultados demonstram que dados históricos
de uma infraestrutura IoT existente podem ser utilizados
para reconstruir temperaturas de sensores físicos com
erros relativamente baixos nas condições experimentais
avaliadas. A comparação entre métodos e configurações
evidencia a importância de considerar as características
de cada sensor-alvo e de selecionar adequadamente as
informações utilizadas na reconstrução, oferecendo uma
base experimental para futuras aplicações de sensoriamento
virtual no LASDPC.